\section{Polynomier}
Polynomier er på formen: \hfill
\(
    p(Z) = a_0 + a_1 Z + a_2 Z^2 + \dots + a_n Z^n
\)

$n \in \mathbb{N}$, $a_i \in \mathbb{R}$, $Z$ en variabel \\

Polynomiets grad er det største $i$ hvor $a_i \neq 0$. \\
Polynomiets ledende koefficient er det tilsvarende $a_i$. \\

Nulpolynomiet har grad $-\infty$.

\begin{align*}
    \deg(p(Z) \cdot q(Z)) &= \deg(p(Z)) + \deg(q(Z)) \\
    \deg(p(Z) + q(Z)) &= \max( \deg(p(Z)), \deg(q(Z)) )
\end{align*}

\subsection{Rødder}

\begin{align*}
    p(\lambda) = 0 &\iff \lambda \text{ er en rod i p} \\
    \lambda \text{ er en rod} &\iff \lambda \text{ deler p så } p(Z) = (Z - \lambda) q(Z) \\
    \deg(p(Z)) = n &\iff p(Z) \text{ har } n \text{ rødder, målt med multiplicitet}
\end{align*}

For komplekse $\lambda$ i polynomier med reelle koefficienter er $\overline{\lambda}$ også en rod.

\kilde{Lemma 5.3.3}

\

Hvis $p(Z) = p_1(Z) \cdot p_2(Z)$ er rødderene af $p(Z)$ rødderne af $p_1(Z)$ og $p_2(Z)$

\kilde{Lemma 5.5.1}


\subsubsection{Rødder af andengradsligninger}
\[
aZ^2 + bZ + c = 0 \hspace{4em} Z = \frac{-b \pm \sqrt{b^2 - 4ac}}{2a}
\]

Hvis diskriminanten $D = b^2 - 4ac$ for andengradsligningen er:
\begin{itemize}
    \item $D > 0$: To forskellige reelle rødder 
    \(
        Z = \frac{-b \pm \sqrt{D}}{2a}
    \)
    \item $D = 0$: Én reel rod
    \(
        Z = \frac{-b}{2a}
    \)
    \item $D < 0$: To komplekse rødder - $\lambda$ og $\overline{\lambda}$
\end{itemize}

\subsection{Multiplicitet og faktorisering}

Et polynomium $p(Z) \in \mathbb{R}[Z]$ med reelle koefficienter kan faktoriseres som produktet af 1. og 2. grads polynomier med reelle koefficienter.

\kilde{Sætning 5.6.3}

Et polynomium $p(Z) \in \mathbb{C}[Z]$ med komplekse koefficienter kan faktoriseres som produktet af 1. grads polynomier med komplekse koefficienter.

\kilde{Korollar 5.6.4}

\subsection{Binomier}
Binomier er på formen \(z^n = w\) for et \(n \in \mathbb{N}\) og et \(w \in \{ \mathbb{C} \setminus 0 \} \) og har $n$ løsninger.

\subsubsection{Polær løsning}
\[
    z = \sqrt[n]{|w|} \cdot e^{i \left( \frac{\arg(w)}{n} + p \cdot \frac{2 \pi}{n} \right)} \quad \text{for } p = 0, 1, \dots, n-1
\]
\subsubsection{Rektangulær løsning}
\[
    z = \sqrt[n]{|w|} \left( \cos\left( \frac{\arg(w)}{n} + p \cdot \frac{2 \pi}{n} \right) + i \sin\left( \frac{\arg(w)}{n} + p \cdot \frac{2 \pi}{n} \right) \right) \quad \text{for } p = 0, 1, \dots, n-1
\]

\subsubsection{Sætrilfælde ved 4. grads binomier}
For et binomie på formen \(Z^4 - w\) kan de resterende rødder findes hvis du kender en af rødderne ved:
\[
    z_2 = i z_1, \quad z_3 = - z_1, \quad z_4 = -i z_1 \kildetag{Eksempel 5.4.1}
\]

Mere generelt kan resterende rødder i et $n$'te grads binomie findes ved at bruge en kendt rod \(z_1\) i formlen:
\[
    z_k = z_1 \cdot e^{i \cdot (k-1) \cdot \frac{2 \pi}{n}} \quad \text{for } k = 1, 2, \dots, n
    \kildetag{\text{Ikke i bogen}}
\]

\subsection{Divisions-Algoritmen}
\eksempel{
Del \(2Z^4 - 6Z^3 + 8Z^2 - 12Z + 8\) med \(Z^2 - 3Z + 2\).

\begin{align*}
    Z^2-3Z+2) \quad 
    2Z^4-6Z^3+8Z^2-12Z+8& \quad (2Z^2+4 \\
    2Z^4-6Z^3+4Z^2 \hspace{4.7em}& \\
    \overline{\hspace{6em} 4Z^2-12Z+8} \\
    4Z^2-12Z+8 \\
    \overline{\hspace{6em} 0}
\end{align*}
Vi får kvotienten \(2Z^2 + 4\) og resten \(0\).
}

Divisions-algoritmen virker ved at man tager sin divisor, og finder hvad den skal ganges med for at udeligne det højeste grad led i dividenden. Dette gøres gentagne gange indtil graden af resten er mindre end graden af divisoren.